\emph{Step 1 (the map on the flag).} The proof of Theorem \claimref{odd-extensions.two-towers-one-curved-universe} gives $\Phi$ and $h'(\Phi X, \Phi Y) = -h(X, Y)$; a linear map with this property sends $h$-isotropic subspaces to $h'$-isotropic ones of the same dimension and keeps the pairing $W \times W' \to \C$ nondegenerate, so $(\Phi W, \Phi W')$ is a Witt pair of $h'$, and $\dim(\Phi W_{1} \cap \Phi W_{2}) = \dim(W_{1} \cap W_{2})$ for any two maximal isotropic subspaces, so the family relation is kept. In a Witt basis $n_{a}, p_{1}, p_{2} \mid n_{b}, q_{1}, q_{2}$ with $n_{a}$ the $\sigma$-direction, $n_{b}$ the $\rho$-direction and $p_{i}, q_{i} \in \mu_{\C}$, $\Phi$ is $\mathrm{diag}(1,1,1,-1,1,1)$, so $\Phi W = W$ and $\Phi W' = W'$ as subspaces.

\emph{Step 2 (the Fock modules).} Let $c(v) = v \wedge$ for $v \in W'$ and $c(w) = 2\iota_{h(w, \cdot)}$ for $w \in W$, so that $c(x)^{2} = h(x,x)$; likewise $c'$ for $h'$. Then $c'(\Phi v)\,\Lambda\Phi = \Lambda\Phi\,c(v)$ and, because $h'(\Phi w, \Phi v) = -h(w, v)$, $c'(\Phi w)\,\Lambda\Phi = -\Lambda\Phi\,c(w)$. The mode number operator is $N_{j} = c(v_{j})c(w_{j})/2h(w_{j}, v_{j})$, and for the mirror $N'_{j} = c'(\Phi v_{j})c'(\Phi w_{j})/2h'(\Phi w_{j}, \Phi v_{j})$; the sign of the product and the sign of the normalisation cancel, so $N'_{j}\,\Lambda\Phi = \Lambda\Phi\,N_{j}$, and the grade $\sum_{j} N_{j}$ is intertwined. The particle-antiparticle labelling is the grade pairing $k \leftrightarrow 3 - k$ of Proposition \claimref{dark-matter.orientation-arrow-seventeenth-refutation}, so it is carried with sign $+1$, and so is any asymmetry measured against the flag. The flag conjugate $C$ exchanges $v_{j}$ and $w_{j}$; it is an isometry of $h$, $\Phi \circ C$ is an anti-isometry carrying $W$ onto $\Phi W'$, and $W$, $W'$ lie in opposite families since they meet in $0$ and $m = 3$ is odd (Theorem \claimref{dark-matter.flag-bit-exists-exactly-tower-dimension}); so $\Phi \circ C$ carries grade $k$ to grade $3 - k$ and the asymmetry with sign $-1$. The check verifies all of this in exact rational arithmetic on the Witt basis, with the signatures $(2,4)$ and $(4,2)$ of the real forms.

\emph{Step 3 (route one needs the orientation).} Clause (v) of Theorem \claimref{dark-matter.mirror-tower-copy-one} equates the two towers' asymmetries along the reversal orbit. Measured against each tower's own flag, the equation reads $\epsilon_{2} = +\epsilon_{6}$ if the mirror's flag is $\Phi W$ and $-\epsilon_{6}$ if it is $\Phi W'$, by Step 2. Applying the flag conjugate to the mirror's Clifford bundle alone is, under the first premise of Proposition \claimref{dark-matter.orientation-arrow-seventeenth-refutation} (conjugation is an isomorphism of every kinematic structure), an automorphism of the two towers' joint kinematic data, the towers sharing only the metric, which it fixes; it exchanges the two assignments. By that proposition's covariance argument, no kinematic or dynamical result fixes the assignment, and it is supplied by a premise: P1's orientation clause, or P2 read as refusing a per-tower orientation, the reading the P1 attack note gives it.

\emph{Step 4 (the reach).} Proposition \claimref{dark-tower-matter.mirrors-temperature-p2} applies P2, at the reach clause (v) uses (a difference of two continuous data of one reversal orbit is an uncaused distinction), to the temperature ratio, a datum of the descent of the same kind, and obtains $\xi = 1$, $\Delta\Neff = 7.45$, excluded. Route one at that reach therefore yields an excluded consequence alongside its intended one; it can be kept only with a scoping of P2 that separates the two data, which P2 at version 2 states (labels, not states; the stored charge is a label and, in the boundary-mode model of Theorem \claimref{synthesis.tower-charge-sector-label-temperature-not}, the temperature is not on any finite window, nor on the full algebra when the Powers-Stormer sum of its (iii) converges), and at that scope route one stands; when infinitely many boundary modes share a frequency band that theorem's scoping reaches $\xi$ and gives $\xi = 1$, and route one does not stand there. Route two does not use P2 at that reach. So under P2 at version 2 clause (v) has both routes.
